Como mencionado na seção anterior, não basta elaborar o algoritmo que responde
à pergunta de pesquisa. É necessário garantir que o algoritmo proposto funcione
em todas as instâncias do problema e apresente desempenho superior ao da
reconstrução completa da árvore de componentes. Nesta seção propõem-se os
passos necessários para assegurar que o algoritmo desenvolvido seja eficiente
e correto em todas as instâncias do problema.

Em primeiro lugar, será pesquisado e construído um arcabouço teórico. Esse
arcabouço fornecerá o rigor matemático necessário para demonstrar a
correção e as propriedades de complexidade do algoritmo proposto. Além da
construção teórica, inclui-se uma etapa dedicada a demonstrações formais do
arcabouço.

Em segundo lugar, com o alicerce matemático estabelecido, será elaborada a
versão algorítmica que realiza a fragmentação ("split") de nós em uma
\textit{max tree}. Caso a construção de um algoritmo eficiente se prove
impossível, o arcabouço permitirá justificar formalmente essa impossibilidade.
Quando o algoritmo for proposto, pretende-se demonstrar sua corretude e
analisar sua complexidade em termos de tempo e espaço. Nesta etapa será
apresentado um pseudocódigo.

Em terceiro lugar, o algoritmo será implementado em uma linguagem de
programação a ser definida (por exemplo, C++ ou Python) para demonstrar a
viabilidade prática da solução. A implementação será utilizada em conjuntos de
dados para coletar métricas que permitam comparar o método proposto com a
reconstrução completa da \textit{max tree}. Antes da execução dos testes, serão
definidas explicitamente as métricas (por exemplo, tempo de execução, uso de
memória, e precisão das propriedades estruturais) e os conjuntos de dados a
serem utilizados.

Por fim, relizar-se-á um estudo de caso para demonstrar, na prática, os
contextos nos quais o operador de \textit{split} pode trazer benefício em
relação às abordagens tradicionais.